\begin{Theorem}\label{th.Gr} An irreducible period matrix $\tau\in \HH_g$ is the period matrix
of a hyperelliptic Jacobian with the basis of cycles chosen in such a way that the corresponding special fundamental system is the one given in~\eqref{eq:FS} if and only if the following cubic identities for second order theta functions are satisfied for all $\s\in \ZZ_2^g$ and for all $z\in \CC^g$:
 \begin{gather}\label{eq:Cu1}
 R_{\s}:= Q[0,0](\tau, z)\T[\s](\tau,z)- \sum_{k=0}^g (-1)^{\langle \s, e_{k+1}\rangle}
Q[s_k, e_{k+1}](\tau, z)\T[\s+s_k](\tau,z)=0,
\end{gather}
where we assume that not all the terms $ Q[0,0](\tau, z)$ and $Q[s_k, e_{k+1}](\tau, z)$ are identically zero in~$z$.
\end{Theorem}
\begin{Remark} The $2^g$ cubic polynomials appearing in (\ref{eq:Cu1}) span an irreducible representation of the Heisenberg group and are of the form
\[R_{\s}=(1,\s,0)R_0.
\]
\end{Remark}

We recall that the result $R_{\s}=0$ is obtained by computing the coefficients in the
formula given in \cite{BK2}.

Indeed, in \cite{Gr} the following relations appear
\begin{gather}
 \sum_{\ep\in\ZZ_2^g} \T[\ep](\tau,z)\T[\ep](\tau,z)\T[\s](\tau,z)\nonumber\\
\qquad {}= \sum_{k=0}^g \sum_{\ep\in\ZZ_2^g} (-1)^{\langle \ep+\s, e_{k+1}\rangle}\T[\ep](\tau,z)\T[\ep+ s_k](\tau,z)\T[\s+s_k](\tau,z). \label{eq:Cu}
\end{gather}

It is an immediate consequence of (\ref{eq:vT}) and the definition of $ Q[\ep,\de](\tau, z)$, that (\ref{eq:Cu1}) and~(\ref{eq:Cu}) are the same equations.

The condition on the non-identically zero terms in Theorem~\ref{th.Gr} turns into the assumption that not all $\theta\left[\begin{smallmatrix}\ep \\ \de\end{smallmatrix}\right](\tau,0) $ are null for
 \[
 \begin{bmatrix} \ep \\ \de\end{bmatrix} =\begin{bmatrix} 0 \\ 0 \end{bmatrix} ,\begin{bmatrix} 0 \\ e_{1}\end{bmatrix}, \dots\begin{bmatrix} s_k \\ e_{k+1} \end{bmatrix},\dots,\begin{bmatrix} s_g \\ 0
\end{bmatrix}.
\]
We observe that this assumption is required in Mumford's characterization of the hyperelliptic locus in \cite{mu}. Actually, the non-vanishing assumption (\ref{nonvanishing}) is stronger.

As proved in \cite{Gr}, these equations in genus 3 are
equivalent to the vanishing of one theta constant, yet when it comes to the genus~4 case they appear to be new; however, as already said, this is not the case, as we will explain. First of all, we need to recall from \cite[Theorem~7.1]{mu}, the so-called generalized Frobenius theta formula. For $U\subset B$ as in (\ref{BU}) we set $\ep_U(j):= \pm1$ according as $j\in U$ or $j\notin U$; then we have

\begin{Theorem} Let $\tau\in\HH_g$ satisfy the vanishing conditions in~\eqref{vanishing}. Then $\forall\, z_1, z_2, z_3, z_4 \in\CC^g$ such that $ \sum\limits_{i=1}^4 z_i=0$:
\begin{gather}\label{eq:FR} \quad \sum_{j\in B} \ep_U(j)\prod_{i=1}^4\theta[m_j](\tau, z_i)=0 . \end{gather}
\end{Theorem}

In \cite{mu} extra variables $a_i$ actually appear, but, as shown in \cite{po}, they are redundant.

We remark that if we apply the addition formula for theta functions to Grushevsky's cubic equations, we obtain similar equations that are also consequence of the generalized Frobenius theta formula; in fact, we have the following:

\begin{Proposition}\label{statement:GF'} If $\tau\in \HH_g$ satisfies equations \eqref{eq:Cu1},
 then, for all $z\in \CC^g$:
\begin{gather}\label{eq:RF2a}
 Q[0, 0] (\tau, z)Q[\ep_1, \de_1 ] (\tau, z)=\sum_{k=0}^g Q[s_k, e_{k+1}] (\tau, z)Q[s_k+\ep_1, e_{k+1}+\de_1] (\tau, z) .
\end{gather}
\end{Proposition}

\begin{proof} We only need to consider
\begin{gather*}
Q[0,0](\tau, z)\sum_{ \s} (-1)^{\langle \s,\de_1\rangle}\T[\s](\tau,z)\T[\s+\ep_1](\tau,z)\\
\qquad{}= \sum_{k=0}^g Q[s_k, e_{k+1}](\tau, z)\sum_{ \s} (-1)^{\langle \s,\de_1+e_{k+1}\rangle}
\T[\s+s_k](\tau,z)\T[\s+\ep_1](\tau,z).
\end{gather*}

Now, the l.h.s.\ yields $ Q[0, 0](\tau, z)Q[\ep_1, \de_1 ] (\tau, z)$. Moreover, we have
\[
Q[s_k+\ep_1, e_{k+1}+\de_1](\tau, z)= \sum_{\s \in\ZZ_2^g} (-1)^{\langle\s, e_{k+1} +\de_1\rangle} \T[\s+s_k+\ep_1](\tau,z)\T[\s](\tau,z).
\]

Hence, in the r.h.s.\ we obtain
\[
\sum_{k=0}^g Q[s_k, e_{k+1}] (\tau, z)Q[s_k+\ep_1, e_{k+1}+\de_1] (\tau, z).\tag*{\qed}
\]\renewcommand{\qed}{}
\end{proof}

\begin{Remark} As a referee correctly pointed out, in the above proposition we are computing $M(\chi)R_0$. \end{Remark}

As a particular yet fundamental case, we will consider $M(0)R_0$, i.e., the case obtained by setting $\ep_1=\de_1=0$; hence we get

\begin{Corollary}\label{statement:GF} If $\tau\in \HH_g$ satisfies equations \eqref{eq:Cu1}, for all $z\in \CC^g$ we have
\begin{gather}\label{eq:RF}
 Q[0, 0]^2(\tau, z)=\sum_{k=0}^g Q[s_k, e_{k+1}]^2(\tau, z).
\end{gather}
\end{Corollary}

\begin{Remark} The important fact is that this formula is exactly the one given in \cite[Corollary~7.5, p.~113]{mu} when $S=\varnothing$, which is a particular case of the general formula obtained as a~consequence of the vanishing conditions~(\ref{vanishing}).
 \end{Remark}

Formula (\ref{eq:RF}) is therefore a special case of Frobenius' theta formula and it is fundamental in Mumford's investigation of Neumann's dynamical system (see \cite[Lemma~9.7]{mu}). It is worth recalling that this formula is obtained by plugging the vanishing conditions in a particular biquadratic Riemann relation of the type (\ref{eq:RR}) with $v_{\ep,\de}=\pm 1$ and $\s=\rho=0$; in fact, we can get it from (\ref{eq:FR}) by setting $z_1=z_2=0$, $z_3=z$, $z_4=-z$ and $ m_1=m_2=m_3=m_4.$

To obtain all of Frobenius' relations described by Mumford in the above cited Corollary, we can evaluate equation (\ref{eq:RF}) at the points $z+(\tau x+y)/4$ for $x,y \in\ZZ^g$, in the same way we did with Riemann's relations to obtain (\ref{eq:RR2}); this leads to the following
\begin{Corollary} If $\tau\in \HH_g$ is a period matrix
satisfying equation \eqref{eq:RF}, then
\begin{gather}\label{eq:RF2}
 \tt{0}{0}(\tau,0)^2\tt{x}{y }(\tau, 2z)^2 =\sum_{k=0}^g \tt{s_k}{e_{k+1}}(\tau,0)^2\tt{s_k+x}{e_{k+1}+y }(\tau, 2z)^2 .
\end{gather}
\end{Corollary}

The converse of Corollary~\ref{statement:GF} also holds:

 \begin{Proposition}\label{statement:FG} Let $\tau\in\HH_g$ satisfy the vanishing conditions in~\eqref{vanishing}; then equation \eqref{eq:RF} implies equations~\eqref{eq:Cu1}.
 \end{Proposition}
 \begin{proof}We know that equation (\ref{eq:RF}) is obtained by evaluating a suitable biquadratic Riemann relation of the form
 \[
\sum \pm Q[\ep, \de]^2 =0.
\]
But this is a trivial relation among the $X_{\s}$, hence, considering the inverse of the map $M(0)$ (cf.\ Proposition~\ref{pr.vG2}) we deduce that the derivative with respect to $X_0$ also gives a trivial relation in the $X_{\s}$ and the same holds for the derivative with respect to any of the $X_{\s}$. Evaluating at the $\T[\s](\tau, z)$ we get the non-trivial relations $R_{\s}=0$.
\end{proof}
